\chapter{Monte Carlo Pricers in Practice}\label{ch:dv:monte-carlo-pricers-in-practice}

A Bermudan option on the best of five shares has no finite-difference grid that fits in memory: a hundred points per
dimension is ten billion nodes. A Monte Carlo engine prices it in seconds, and gets the exercise decision wrong by
design. It estimates the value of waiting from a regression on the simulated paths, which is only an approximation of
the true continuation value, so the exercise policy it follows is suboptimal and its price is a lower bound. A good desk
therefore reports two numbers, a lower bound from the policy and an upper bound from a dual method, and the gap between
them measures how much the regression still misses. This chapter builds the machinery around the Monte Carlo method of
One Quant Book 4: simulating the models, early exercise by regression and its dual bound, three ways to compute \omterm{def:dv:greeks-and-the-hedging-pnl:greeks}{Greeks},
low-discrepancy numbers with the Brownian bridge, and the engineering that turns it into a pricer.

\section{Simulating the models}

Lognormal assets are simulated exactly, date to date. Every other model is discretised, and the error of the scheme is a
bias that adds to the statistical error and does not average away. Heston's variance is the standard hard case. An Euler step
can make the variance negative, and truncating it at zero leaves a bias that shrinks only slowly with the step size. Andersen
designed a scheme that matches the first two moments of the variance's exact transition.

\begin{definition}[Quadratic-exponential scheme]\label{def:dv:monte-carlo-pricers-in-practice:qe}
The \emph{quadratic-exponential scheme}\index{quadratic-exponential scheme} (QE) simulates Heston's variance over a step by matching
the mean $m$ and variance $s^2$ of its exact conditional law: when $\psi=s^2/m^2$ is small it draws $a(b+Z)^2$, a scaled
non-central square of a normal; when $\psi$ is large, a mixture of a mass at zero and an exponential. The log-price then uses the
integrated variance implied by the start and end variances.
\end{definition}

\begin{example}[Euler against QE]\label{ex:dv:monte-carlo-pricers-in-practice:qe}
A one-year at-the-money call in chapter 10's base \omterm{def:dv:stochastic-volatility:heston}{Heston model} ($v_0=\bar v=0.04$, $\kappa=1.5$, $\eta=0.6$, $\rho=-0.7$) is worth 6.751 by
Fourier inversion. With 400\,000 paths and full-truncation Euler, the bias is 0.857 with 4 steps a year, 0.373 with 8, 0.140 with 16 and
0.041 with 32. With QE the price is within one standard error (0.014) of the exact value at every step count, from 4 steps a year
(\cref{fig:dv:monte-carlo-pricers-in-practice:heston}). QE saves a factor of ten or more in steps, and so in run time.
\end{example}

\begin{omfigure}
\begin{tikzpicture}
\begin{axis}[width=10.5cm,height=5cm,xmode=log,ymode=log,log basis x=2,
  xlabel={time steps a year},ylabel={absolute bias},
  tick label style={font=\small},label style={font=\small},
  xmin=3,xmax=40,ymin=1e-4,ymax=2,grid=major,grid style={gray!25},
  xtick={4,8,16,32},xticklabels={4,8,16,32},
  legend columns=3,legend cell align=left,legend style={font=\footnotesize,at={(0.5,-0.3)},anchor=north,draw=none,fill=none,/tikz/every even column/.append style={column sep=8pt}},
  table/col sep=comma]
\addplot[omProp,very thick,mark=*,mark size=1.6pt] table[x=steps,y=euler]{figdata/derivatives/23-monte-carlo-pricers-in-practice/heston_bias.csv};
\addplot[omThm,very thick,mark=square*,mark size=1.6pt] table[x=steps,y=qe]{figdata/derivatives/23-monte-carlo-pricers-in-practice/heston_bias.csv};
\addplot[gray,thick,dashed] table[x=steps,y=se]{figdata/derivatives/23-monte-carlo-pricers-in-practice/heston_bias.csv};
\legend{full-truncation Euler,quadratic-exponential,one standard error}
\end{axis}
\end{tikzpicture}

\omcaption{Bias of a one-year at-the-money call in a \omterm{def:dv:stochastic-volatility:heston}{Heston model} whose variance can reach zero, by time steps, against the Fourier price
(400\,000 paths). Euler's bias decays slowly; QE's stays within the statistical error at every step count. Data: the
tutorial.}\label{fig:dv:monte-carlo-pricers-in-practice:heston}
\end{omfigure}

\section{Early exercise by regression}

Chapter 6 priced Bermudan and American options backwards on a tree, where the continuation value at each node is an average over its
children. A simulation moves forward and has no such average. Longstaff and Schwartz estimated it by regression.

\begin{definition}[Longstaff--Schwartz method]\label{def:dv:monte-carlo-pricers-in-practice:lsm}
The \emph{Longstaff--Schwartz method}\index{Longstaff--Schwartz method} prices a Bermudan option by simulation: working backwards over the
exercise dates, it regresses the discounted cash flows that each path realises from the next date on, against a set of basis functions of the
state, using the paths that are in the money, and exercises where the exercise value exceeds the fitted continuation value; the policy is then
applied to independent paths to give a lower bound.
\end{definition}

The price from an estimated policy is biased low, because no policy beats the optimal one. Fitting and pricing on the same paths adds an upward
bias, because the regression has seen the paths it then prices. The build fits on one set of paths and prices on another, which keeps the result
a lower bound. How low it is depends on the basis. A regression that cannot represent the continuation value leads to exercising at the wrong
times.

\begin{definition}[Dual upper bound]\label{def:dv:monte-carlo-pricers-in-practice:dual}
The \emph{dual upper bound}\index{dual upper bound} of a Bermudan option is $\E\bigl[\max_j(h_j-M_j)\bigr]$ for a martingale $M$ with $M_0=0$
and discounted exercise values $h_j$; it exceeds the price for every martingale and equals it for the martingale part of the value process. The
Andersen--Broadie algorithm builds $M$ from a lower-bound policy by nested simulation, so that lower and upper bounds together give an interval
for the price.
\end{definition}

In the build, the policy's continuation value $Q_j$ is estimated at every outer node by inner paths that follow the policy. With $L_j$ the
discounted exercise value where the policy exercises and $Q_j$ where it continues, the martingale's increments are $L_j-Q_{j-1}$. The inner
estimates are noisy, and the maximum in the bound turns noise into an upward bias. The bound is valid for any number of inner paths, but tight only
for many.

\begin{example}[The bias of too few inner paths]\label{ex:dv:monte-carlo-pricers-in-practice:inner}
A one-year Bermudan put with ten exercise dates (spot and strike 100, 5\%, 20\%) is worth 6.0326 on chapter 6's tree. Longstaff--Schwartz with a
cubic basis gives 6.0325 (standard error 0.023). The dual bound on 1\,000 outer paths is 6.134 with 100 inner paths per node, 6.054 with 400 and
6.037 with 1\,600 (\cref{fig:dv:monte-carlo-pricers-in-practice:bounds}, right).
\end{example}

\begin{example}[A three-asset Bermudan max-call]\label{ex:dv:monte-carlo-pricers-in-practice:maxcall}
A call on the best of three independent shares (each at 100, strike 100, volatility 20\%, dividend yield 10\%, rate 5\%), exercisable at nine dates
over three years. The basis for the regression is a set of polynomials in the two largest prices, their product, the product of all three and the
intrinsic value. With it, Longstaff--Schwartz on 200\,000 fresh paths gives 18.649 (standard error 0.039), and the dual bound (1\,000 outer paths,
1\,000 inner) gives 18.713 (0.019). The gap is 0.064, 0.3\% of the price. With a basis of only the largest price and its square, the lower bound
falls to 18.137 and the upper rises to 18.805, a gap of 0.668, ten times wider
(\cref{fig:dv:monte-carlo-pricers-in-practice:bounds}, left).
\end{example}

\begin{omfigure}
\begin{tikzpicture}
\begin{groupplot}[group style={group size=2 by 1,horizontal sep=1.6cm},
  width=6.6cm,height=5cm,
  tick label style={font=\small},label style={font=\small},grid=major,grid style={gray!25},table/col sep=comma]
\nextgroupplot[xlabel={regression basis},ylabel={price},xmin=-0.6,xmax=1.6,ymin=18.0,ymax=18.95,
  xtick={0,1},xticklabels={rich,small},title={\small max-call bounds},
  legend columns=1,legend cell align=left,legend style={font=\footnotesize,at={(0.5,-0.32)},anchor=north,draw=none,fill=none}]
\addplot[omDef,only marks,mark=square*,mark size=2.4pt,error bars/.cd,y dir=both,y explicit] table[x=i,y=lower,y error=lower_se]{figdata/derivatives/23-monte-carlo-pricers-in-practice/bounds.csv};
\addplot[omThm,only marks,mark=triangle*,mark size=2.8pt,error bars/.cd,y dir=both,y explicit] table[x=i,y=upper,y error=upper_se]{figdata/derivatives/23-monte-carlo-pricers-in-practice/bounds.csv};
\legend{Longstaff--Schwartz (lower),dual (upper)}
\nextgroupplot[xmode=log,log basis x=2,xlabel={inner paths per node},ylabel={dual upper bound},xmin=70,xmax=2200,ymin=6.0,ymax=6.16,
  xtick={100,400,1600},xticklabels={100,400,1\,600},
  yticklabel style={/pgf/number format/fixed,/pgf/number format/precision=2},title={\small Bermudan put},
  legend columns=1,legend cell align=left,legend style={font=\footnotesize,at={(0.5,-0.32)},anchor=north,draw=none,fill=none}]
\addplot[omThm,very thick,mark=*,mark size=1.6pt] table[x=inner,y=upper]{figdata/derivatives/23-monte-carlo-pricers-in-practice/inner_bias.csv};
\addplot[black,dashed] table[x=inner,y=tree]{figdata/derivatives/23-monte-carlo-pricers-in-practice/inner_bias.csv};
\legend{dual bound,tree value}
\end{groupplot}
\end{tikzpicture}

\omcaption{Left: lower and upper bounds (bars of one standard error) for a three-asset Bermudan max-call with a rich and a small regression basis: the
gap is the price of a poor basis. Right: the dual bound for a one-asset Bermudan put against the number of inner paths per node, approaching the
tree's value from above. Data: the tutorial.}\label{fig:dv:monte-carlo-pricers-in-practice:bounds}
\end{omfigure}

\section{Greeks: pathwise, likelihood ratio, adjoint}

Bumping a Monte Carlo price and repricing gives \omterm{def:dv:greeks-and-the-hedging-pnl:greeks}{Greeks} with the variance of a difference of two noisy numbers. Common random numbers make the
difference far less noisy. Two direct methods, due to Broadie and Glasserman, avoid the bump altogether and give \omterm{def:dv:greeks-and-the-hedging-pnl:greeks}{Greeks} from the same paths as the
price.

\begin{definition}[Pathwise Greek]\label{def:dv:monte-carlo-pricers-in-practice:pathwise}
A \emph{pathwise Greek}\index{pathwise Greek} differentiates the discounted payoff along each simulated path with respect to the parameter, holding
the random numbers fixed, and averages; it requires the payoff to be continuous in the parameter (almost surely differentiable), and fails for
digitals and barriers.
\end{definition}

\begin{definition}[Likelihood-ratio Greek]\label{def:dv:monte-carlo-pricers-in-practice:lr}
A \emph{likelihood-ratio Greek}\index{likelihood-ratio Greek} differentiates the density of the simulated state instead of the payoff: it averages the
discounted payoff times the score, the derivative of the log-density with respect to the parameter; it works for discontinuous payoffs, at the cost
of a higher variance.
\end{definition}

The table compares the two methods on one-year at-the-money options (spot and strike 100, rate 5\%, volatility 20\%) with 200\,000 paths, standard
errors in brackets. The last column applies the pathwise method to the digital smoothed into a call spread.

\begin{center}\small
\begin{tabular}{lrrrr}
\toprule
one year & exact & pathwise & likelihood ratio & smoothed\\
\midrule
call delta & 0.6368 & 0.6384 (0.0013) & 0.6413 (0.0033) & ---\\
call \omterm{def:dv:greeks-and-the-hedging-pnl:greeks}{vega} & 37.52 & 37.72 (0.17) & 38.24 (0.62) & ---\\
digital delta & 0.01876 & 0 (wrong) & 0.01883 (0.00006) & 0.01895 (0.00021)\\
digital \omterm{def:dv:greeks-and-the-hedging-pnl:greeks}{vega} & $-0.657$ & 0 (wrong) & $-0.648$ (0.010) & $-0.663$ (0.007)\\
\bottomrule
\end{tabular}
\end{center}

For the call the pathwise estimator's variance is 6.5 times lower for delta and 13 times lower for \omterm{def:dv:greeks-and-the-hedging-pnl:greeks}{vega}. For the digital it returns zero, since the
payoff's derivative is zero almost everywhere. The likelihood ratio gets it right. Smoothing the digital into a call spread of half-width one restores
a pathwise estimate, at the cost of a bias of the spread's own.

The pathwise method has an efficient implementation that matters most for large books. Giles and Glasserman's adjoint method
runs the same pathwise calculation backwards along each path: it gives the same values, and its advantage grows with the number of
inputs, such as the points of a \omterm{def:dv:implied-volatility-and-its-surface:surface}{volatility surface}.

\begin{example}[Twelve vegas for the price of three passes]\label{ex:dv:monte-carlo-pricers-in-practice:adjoint}
An arithmetic Asian call (one year, 48 steps, strike 100, 20\%) is priced along with its sensitivity to the volatility in each of twelve monthly
buckets. Forward mode would need one pass per bucket. The adjoint (reverse) pass of One Quant Book 4, chapter 28, runs once backwards along each stored
path and accumulates the derivative of the payoff with respect to every step's state and volatility. It returns all twelve \omterm{def:dv:greeks-and-the-hedging-pnl:greeks}{vegas} and the delta at a
cost of a few forward passes. The price is 5.901 and the delta 0.593. The \omterm{def:dv:greeks-and-the-hedging-pnl:greeks}{vega} of the first month is 5.064, where bumping that month's volatility with
the same draws gives 5.065, and the \omterm{def:dv:greeks-and-the-hedging-pnl:greeks}{vegas} fall to 0.016 for the last month (\cref{fig:dv:monte-carlo-pricers-in-practice:adjoint}).
\end{example}

\begin{omfigure}
\begin{tikzpicture}
\begin{axis}[width=10.5cm,height=4.6cm,ybar,bar width=12pt,
  xlabel={volatility bucket (month)},ylabel={vega (per unit of volatility)},
  tick label style={font=\small},label style={font=\small},
  xmin=0.3,xmax=12.7,ymin=0,ymax=5.6,grid=major,grid style={gray!25},xtick={1,2,3,4,5,6,7,8,9,10,11,12},
  table/col sep=comma]
\addplot[omThm,fill=omThm!40] table[x=month,y=vega]{figdata/derivatives/23-monte-carlo-pricers-in-practice/bucket_vega.csv};
\end{axis}
\end{tikzpicture}

\omcaption{Bucketed \omterm{def:dv:greeks-and-the-hedging-pnl:greeks}{vegas} of a one-year arithmetic Asian call from one adjoint pass along each path: the average's sensitivity to early volatility
is largest because early moves carry through every later fixing. Data: the tutorial.}\label{fig:dv:monte-carlo-pricers-in-practice:adjoint}
\end{omfigure}

\section{Bridges and quasi-random numbers}

Low-discrepancy sequences (Book 4, chapter 26) fill the unit cube more evenly than random numbers, and for smooth integrands their error falls faster
than $1/\sqrt N$. A path of $d$ dates is a point in $d$ dimensions, however, and low-discrepancy sequences lose their advantage in the later
coordinates. The Brownian bridge reorders the construction so that the first coordinates carry most of the path's variance. The first normal sets
the endpoint $W_T$, the second the midpoint given the ends, and so on. The good early coordinates of the sequence then determine the path's shape.

\begin{example}[An Asian option with Halton points]\label{ex:dv:monte-carlo-pricers-in-practice:qmc}
A 16-date Asian call (one year, strike 100, 20\%). With 4\,096 paths the Monte Carlo standard error is 0.131. Halton points, randomly shifted twenty times
to measure their error, give 0.027 with the usual step-by-step construction and 0.0092 with the Brownian bridge, fourteen times smaller than Monte Carlo,
the accuracy of about 200 times more random paths. At 256 paths the three errors are 0.53, 0.34 and 0.13
(\cref{fig:dv:monte-carlo-pricers-in-practice:qmc}).
\end{example}

\begin{omfigure}
\begin{tikzpicture}
\begin{axis}[width=10.5cm,height=5cm,xmode=log,ymode=log,log basis x=2,
  xlabel={paths},ylabel={error (standard deviation)},
  tick label style={font=\small},label style={font=\small},
  xmin=200,xmax=5200,ymin=0.005,ymax=0.8,grid=major,grid style={gray!25},
  xtick={256,1024,4096},xticklabels={256,1\,024,4\,096},
  legend columns=3,legend cell align=left,legend style={font=\footnotesize,at={(0.5,-0.3)},anchor=north,draw=none,fill=none,/tikz/every even column/.append style={column sep=8pt}},
  table/col sep=comma]
\addplot[omProp,very thick,mark=*,mark size=1.6pt] table[x=n,y=mc]{figdata/derivatives/23-monte-carlo-pricers-in-practice/qmc.csv};
\addplot[omDef,very thick,mark=square*,mark size=1.6pt] table[x=n,y=halton]{figdata/derivatives/23-monte-carlo-pricers-in-practice/qmc.csv};
\addplot[omThm,very thick,mark=triangle*,mark size=1.8pt] table[x=n,y=bridge]{figdata/derivatives/23-monte-carlo-pricers-in-practice/qmc.csv};
\legend{Monte Carlo,Halton,Halton with bridge}
\end{axis}
\end{tikzpicture}

\omcaption{Error of a 16-date Asian call against the number of paths: Monte Carlo's standard error, and the spread over random shifts of Halton points
with the step-by-step and the Brownian-bridge constructions. Data: the tutorial.}\label{fig:dv:monte-carlo-pricers-in-practice:qmc}
\end{omfigure}

\section{Engineering a Monte Carlo pricer}

A production engine separates what varies from what is shared. Path generators are one per model: lognormal, \omterm{def:dv:local-volatility:model}{local volatility}, Heston-QE,
stochastic-local. They write paths on a date grid that the product needs, and they never know the payoff. Payoffs read paths only (chapters 16 and
18). Estimators, Longstaff--Schwartz, the \omterm{def:dv:greeks-and-the-hedging-pnl:greeks}{Greeks} and the dual bound, sit on top. Four rules keep it honest.
\begin{itemize}
\item \textbf{Seeds are data.} Every price is reproducible from its seed. \omterm{def:dv:greeks-and-the-hedging-pnl:greeks}{Greeks} by bumping use the same draws (common random numbers).
\item \textbf{Every number carries its error.} The standard error for statistics, a convergence check for the discretisation, and a bound gap for early
exercise.
\item \textbf{Variance reduction first.} Antithetics, control variates (chapter 16's geometric Asian), conditioning (chapter 10's \omterm{prop:dv:stochastic-volatility:mixing}{mixing formula}) and
the bridge usually buy more than extra paths.
\item \textbf{Vectorise over paths, loop over time.} The build's generators are arrays of paths advanced date by date, the natural shape for numpy, for
a GPU, and for the adjoint pass that runs back over the same dates.
\end{itemize}

\section{Tutorial: bounds, Greeks and bridges}\label{tut:dv:monte-carlo-pricers-in-practice:tut}

\begin{tutorial}
\textbf{Goal.} Price a Bermudan max-call by Longstaff--Schwartz and bound it above by the dual; compute \omterm{def:dv:greeks-and-the-hedging-pnl:greeks}{Greeks} of a call and a digital by the pathwise
and likelihood-ratio methods and bucketed \omterm{def:dv:greeks-and-the-hedging-pnl:greeks}{vegas} by an adjoint pass; compare Euler and QE for Heston; measure the effect of the Brownian bridge on Halton
points. \textbf{End state:} the four figures, the table and the numbers of the weekend problem.
\begin{tutsteps}
\item \textbf{The regression}, backwards over the exercise dates, in-the-money paths only:
\omcode{code/firm/mcpricer/firm_mcpricer.py}{85}{103}{Longstaff--Schwartz regression.}\label{lst:dv:monte-carlo-pricers-in-practice:lsm}
\item \textbf{The dual bound}: inner continuation values, martingale increments $L_j-Q_{j-1}$, the maximum:
\omcode{code/firm/mcpricer/firm_mcpricer.py}{126}{157}{The Andersen--Broadie upper bound.}\label{lst:dv:monte-carlo-pricers-in-practice:dual}
\item \textbf{The adjoint pass} for bucketed \omterm{def:dv:greeks-and-the-hedging-pnl:greeks}{vegas}:
\omcode{code/firm/mcpricer/firm_mcpricer.py}{184}{206}{Forward and reverse passes for an Asian call.}\label{lst:dv:monte-carlo-pricers-in-practice:adjoint}
\item \textbf{Run} \texttt{dv\_mc.bounds()}, \texttt{put\_bias()}, \texttt{greek\_table()}, \texttt{adjoint()}, \texttt{heston\_bias()}, \texttt{qmc\_study()}
and \texttt{fig\_mc.py}.
\end{tutsteps}
\textbf{What to change next.} Price the max-call with correlation 0.5; add the geometric Asian as a control variate on top of the Halton points; replace the
digital's likelihood ratio by a conditional expectation over the last step (smoothing by conditioning).
\end{tutorial}

\section{Build: the Monte Carlo engine}\label{bld:dv:monte-carlo-pricers-in-practice:mcpricer}

\begin{build}
\bfield{Purpose} The miniature firm's simulation engine: path generators, early exercise with bounds, \omterm{def:dv:greeks-and-the-hedging-pnl:greeks}{Greeks} by three methods, quasi-random numbers;
chapter 28's Monte Carlo engine wraps it.
\bfield{Interface} \texttt{gbm\_paths(s0, vols, corr, times, r, q, n, seed)}; \texttt{heston\_terminal(\ldots, scheme)}; \texttt{lsm\_fit(paths, ex\_idx, payoff,
basis, k, disc)}, \texttt{lsm\_price}, \texttt{exercise\_now}, \texttt{max\_call\_basis}, \texttt{small\_basis}; \texttt{dual\_upper(outer, \ldots, step\_sim,
n\_inner, seed)}; \texttt{call\_greeks(\ldots, digital, smooth)}; \texttt{asian\_adjoint}; \texttt{halton(n, dim, shift)}, \texttt{brownian\_bridge(z, t)},
\texttt{norm\_inv}.
\bfield{Rules} Regression and pricing on independent paths; a lower bound is reported with its \omterm{def:dv:monte-carlo-pricers-in-practice:dual}{dual upper bound} for every Bermudan; \omterm{def:dv:greeks-and-the-hedging-pnl:greeks}{Greeks} by pathwise where
the payoff is continuous, by likelihood ratio or smoothing where it is not; seeds explicit; no memorised constants in the numerics (the inverse normal is
Newton's method on the error function).
\bfield{Acceptance tests} \texttt{code/firm/mcpricer/tests/}: the Bermudan bounds bracket the tree value and the dual bound falls with inner paths; pathwise and
\omterm{def:dv:monte-carlo-pricers-in-practice:lr}{likelihood-ratio Greeks} agree with Black--Scholes and the pathwise digital is zero; the adjoint matches bumping; QE matches the Fourier price; Halton's first
points, the inverse normal and the bridge's covariances.
\bfield{Stretch} Sobol points with published direction numbers (cited, not recalled); multilevel Monte Carlo (Book 4); an adjoint pass through
Longstaff--Schwartz.
\end{build}

\begin{omsources}
\item F.~A.~Longstaff and E.~S.~Schwartz, ``Valuing American options by simulation: a simple least-squares approach'', \emph{Review of Financial Studies}
14(1) (2001) 113--147.
\item L.~Andersen and M.~Broadie, ``Primal-dual simulation algorithm for pricing multidimensional American options'', \emph{Management Science} 50(9)
(2004) 1222--1234.
\item L.~Andersen, ``Simple and efficient simulation of the Heston stochastic volatility model'', \emph{Journal of Computational Finance} 11(3) (2008) 1--42.
\item M.~Broadie and P.~Glasserman, ``Estimating security price derivatives using simulation'', \emph{Management Science} 42(2) (1996) 269--285.
\item P.~Glasserman, \emph{Monte Carlo Methods in Financial Engineering} (Springer, 2003).
\item M.~Giles and P.~Glasserman, ``Smoking adjoints: fast Monte Carlo Greeks'', \emph{Risk} 19 (2006) 88--92; Oxford report NA-05-15 (2005).
\end{omsources}

\section{Exercises}

\begin{exercise}[$\star$]\label{exo:dv:monte-carlo-pricers-in-practice:1}
Why is a Longstaff--Schwartz price, computed on paths independent of the regression, a lower bound?
\end{exercise}

\begin{exercise}[$\star$]\label{exo:dv:monte-carlo-pricers-in-practice:2}
Derive the pathwise delta of a call under Black--Scholes, $e^{-rT}\mathbf 1\{S_T>K\}S_T/S_0$.
\end{exercise}

\begin{exercise}[$\star$]\label{exo:dv:monte-carlo-pricers-in-practice:3}
Why is the pathwise delta of a digital zero, and why is the likelihood-ratio delta not?
\end{exercise}

\begin{exercise}[$\star\star$]\label{exo:dv:monte-carlo-pricers-in-practice:4}
Derive the likelihood-ratio \omterm{def:dv:greeks-and-the-hedging-pnl:greeks}{vega} weight $(Z^2-1)/\sigma-Z\sqrt T$ for a lognormal terminal price.
\end{exercise}

\begin{exercise}[$\star\star$]\label{exo:dv:monte-carlo-pricers-in-practice:5}
Explain why too few inner paths bias the dual bound upwards, never downwards.
\end{exercise}

\begin{exercise}[$\star\star$]\label{exo:dv:monte-carlo-pricers-in-practice:6}
Why does the Brownian bridge help quasi-random numbers but not pseudo-random ones?
\end{exercise}

\begin{exercise}[$\star\star\star$]\label{exo:dv:monte-carlo-pricers-in-practice:7}
\emph{Coding.} Fit and price the max-call's Longstaff--Schwartz policy on the same paths and compare with the independent-paths price. Explain the sign of the
difference.
\end{exercise}

\begin{exercise}[$\star\star\star$]\label{exo:dv:monte-carlo-pricers-in-practice:8}
\emph{Find the flaw.} ``Our Monte Carlo prices the Bermudan to two cents: the standard error is 0.01.''
\end{exercise}

\section{Problem: The Bermudan Bounds}

\begin{problem}[{Weekend problem --- how good is the exercise policy}]\label{pb:dv:monte-carlo-pricers-in-practice:1}
A desk sells a three-year Bermudan call on the best of three shares (each at 100, strike 100, 20\% volatility, 10\% dividend yield, rate 5\%, nine exercise dates)
and prices it by Longstaff--Schwartz.

\medskip\noindent\textbf{Part I --- The lower bound.}
\begin{enumerate}
\item Give the price with the rich basis and its standard error.
\item Why are fitting and pricing done on separate paths?
\item Give the price with the small basis (the largest price and its square).
\item Which exercise decisions does the small basis get wrong?
\item What would a European max-call's price tell you?
\end{enumerate}

\medskip\noindent\textbf{Part II --- The upper bound.}
\begin{enumerate}[resume]
\item Give the dual bound for the rich basis with 1\,000 outer and 1\,000 inner paths, and its standard error.
\item Give the gap and express it as a share of the price.
\item Give the bound and gap for the small basis.
\item On the one-asset Bermudan put, how does the bound change with 100, 400 and 1\,600 inner paths?
\item What does the dual bound cost relative to the lower bound, and why?
\end{enumerate}

\medskip\noindent\textbf{Part III --- Around the price.}
\begin{enumerate}[resume]
\item How would you compute the Bermudan's delta?
\item Why not bump and reprice with fresh paths?
\item What does the adjoint method offer for a book of such trades?
\item How would quasi-random numbers help, and what limits them here?
\item Which discretisation error does this product not have, and which model would introduce one?
\end{enumerate}

\medskip\noindent\textbf{Part IV --- Judgement.}
\begin{enumerate}[resume]
\item Which price would you quote, and which would you book?
\item How much reserve would you hold for the exercise policy?
\item When would you stop improving the basis?
\item State the \emph{named result}: the gap between the Longstaff--Schwartz lower bound and the \omterm{def:dv:monte-carlo-pricers-in-practice:dual}{dual upper bound} for the three-asset Bermudan max-call, and its
sensitivity to the regression basis.
\item In one sentence: what does the gap between the bounds measure?
\end{enumerate}
\end{problem}

\section{Interview questions}

\begin{interviewq}[$\star$ \iqroles{developer, researcher}]\label{iq:dv:monte-carlo-pricers-in-practice:1}
Explain the Longstaff--Schwartz algorithm.
\end{interviewq}

\begin{interviewq}[$\star\star$ \iqroles{researcher}]\label{iq:dv:monte-carlo-pricers-in-practice:2}
How do you get an upper bound for a Bermudan price by simulation?
\end{interviewq}

\begin{interviewq}[$\star\star$ \iqroles{researcher}]\label{iq:dv:monte-carlo-pricers-in-practice:3}
Compare pathwise and \omterm{def:dv:monte-carlo-pricers-in-practice:lr}{likelihood-ratio Greeks}. When does each fail?
\end{interviewq}

\begin{interviewq}[$\star\star$ \iqroles{developer}]\label{iq:dv:monte-carlo-pricers-in-practice:4}
How would you simulate the \omterm{def:dv:stochastic-volatility:heston}{Heston model} efficiently, and why not Euler?
\end{interviewq}

\begin{interviewq}[$\star\star$ \iqroles{developer, risk}]\label{iq:dv:monte-carlo-pricers-in-practice:5}
Your Monte Carlo risk is noisy from day to day. What do you change?
\end{interviewq}

\begin{interviewq}[$\star\star\star$ \iqroles{developer}]\label{iq:dv:monte-carlo-pricers-in-practice:6}
Explain adjoint algorithmic differentiation for Monte Carlo \omterm{def:dv:greeks-and-the-hedging-pnl:greeks}{Greeks}. What does it cost in memory and time?
\end{interviewq}
